\hypertarget{ux5deux5e0ux5d2ux5e0ux5d5ux5e0ux5d9-ux5e7ux5e9ux5d1}{%
\section{מנגנוני
קשב}\label{ux5deux5e0ux5d2ux5e0ux5d5ux5e0ux5d9-ux5e7ux5e9ux5d1}}

\hypertarget{ux5e7ux5e9ux5d1-ux5e8ux5da-ux5d5ux5e6ux5d9ux5d5ux5e0ux5d9ux5dd-ux5e9ux5dc-ux5d9ux5d9ux5e9ux5d5ux5e8}{%
\subsection{קשב רך וציונים של
יישור}\label{ux5e7ux5e9ux5d1-ux5e8ux5da-ux5d5ux5e6ux5d9ux5d5ux5e0ux5d9ux5dd-ux5e9ux5dc-ux5d9ux5d9ux5e9ux5d5ux5e8}}

מנגנון הקשב הרך, כפי שהוצג על ידי Bahdanau et al.~(2015), שינה באופן
מהותי את הגישה לבעיות תרגום מכונה רצף-אל-רצף. עד להצגת מנגנון זה, מודלי
קידוד-פענוח קלאסיים נשענו על וקטור הקשר הסופי של המקודד כסיכום יחיד של
רצף הקלט כולו --- גישה המגבילה את יכולת המודל לטפל ברצפים ארוכים. הקשב
הרך מתגבר על מגבלה זו בכך שהוא מאפשר לפענוחן להסתכל על כלל מצבי המקודד
בכל שלב של הפקת הפלט, ולהקצות משקל שונה לכל מצב בהתאם לרלוונטיות שלו
לצעד הנוכחי.

באופן פורמלי, נסמן את מצבי המקודד הנסתרים בתור h\_1, \ldots, h\_T, ואת
מצב הפענוחן בצעד t בתור s\_t. ציון היישור בין מצב הפענוחן s\_t לכל מצב
מקודד h\_j מחושב באמצעות פונקציית ציון, המוגדרת לרוב כשכבת רשת עצבית
קטנה:

e\_\{t,j\} = a(s\_\{t-1\}, h\_j) = v\_a\^{}T * tanh(W\_a * s\_\{t-1\} +
U\_a * h\_j)

כאשר W\_a, U\_a, v\_a הם פרמטרים הניתנים ללמידה. ציוני היישור עוברים
נרמול באמצעות פונקציית softmax כדי לייצר משקלי קשב:

alpha\_\{t,j\} = softmax(e\_\{t,j\}) = exp(e\_\{t,j\}) /
sum\_k(exp(e\_\{t,k\}))

וקטור ההקשר c\_t המוזן לפענוחן בכל צעד t הוא הממוצע המשוקלל של מצבי
המקודד:

c\_t = sum\_j(alpha\_\{t,j\} * h\_j)

מנגנון זה הניב שיפורים ניכרים באיכות תרגום מכונה נוירוני, ובמיוחד בתרגום
משפטים ארוכים. יתרה מזאת, משקלי הקשב ניתנים לפרשנות: הם מגלים אילו מילים
בשפת המקור המודל ``השקיף'' עליהן בעת הפקת כל מילה בשפת היעד, ובכך מספקים
רמת שקיפות נדירה בהבנת פעולת הרשת העצבית.

\hypertarget{ux5e7ux5e9ux5d1-ux5e2ux5e6ux5deux5d9-ux5d5ux5deux5dbux5e4ux5dcux5d4-ux5e4ux5e0ux5d9ux5deux5d9ux5ea-ux5deux5d3ux5d5ux5e8ux5d2ux5ea}{%
\subsection{קשב-עצמי ומכפלה פנימית
מדורגת}\label{ux5e7ux5e9ux5d1-ux5e2ux5e6ux5deux5d9-ux5d5ux5deux5dbux5e4ux5dcux5d4-ux5e4ux5e0ux5d9ux5deux5d9ux5ea-ux5deux5d3ux5d5ux5e8ux5d2ux5ea}}

מנגנון הקשב העצמי, שהוצג במסגרת ארכיטקטורת ה-Transformer על ידי Vaswani
et al.~(2017), מרחיב את רעיון הקשב לסיטואציה שבה הרצף הוא מקור לעצמו.
במקום שפענוחן יסתכל על מצבי מקודד חיצוניים, כל אלמנט ברצף בוחן את קשריו
עם כל האלמנטים האחרים באותו רצף --- מה שמאפשר ללמוד ייצוגים עשירים
בהקשר.

בלב מנגנון הקשב העצמי עומדות שלוש מטריצות: שאילתה (Q --- Query), מפתח (K
--- Key) וערך (V --- Value). עבור רצף קלט X, שלוש המטריצות מחושבות על
ידי הטלות לינאריות:

Q = X * W\^{}Q, K = X * W\^{}K, V = X * W\^{}V

כאשר W\^{}Q, W\^{}K, W\^{}V הן מטריצות פרמטרים הניתנות ללמידה ממימד
d\_model x d\_k. פונקציית הקשב מחושבת לאחר מכן כ:

Attention(Q, K, V) = softmax(Q * K\^{}T / sqrt(d\_k)) * V

החלוקה בשורש d\_k היא מרכיב קריטי: כאשר ממדיות d\_k גדולה, מכפלות הנקודה
בין וקטורי השאילתה והמפתח עלולות לגדול לערכים גדולים, מה שדוחף את
פונקציית ה-softmax לאזורים עם שיפועים קטנים מאוד (רוויה), ובכך מקשה על
האימון. החלוקה בשורש d\_k מייצבת את השיפועים ומסייעת להתכנסות.

ניתן לפרש את פעולת מנגנון הקשב העצמי כחיפוש עם-ממשק-רך במאגר מידע: לכל
אסימון (שאילתה), המנגנון משווה את ``שאלתו'' מול ``מפתחות'' כל האסימונים
האחרים, מחשב ציוני קרבה, ולאחר מכן מחזיר שקלול של ה''ערכים'' של כל
האסימונים על פי ציונים אלו. בשונה מרשתות נשנות (RNN), פעולה זו מבוצעת על
כל האסימונים באופן מקבילי, מה שמפחית דרסטית את זמני האימון ומאפשר למידת
תלויות ארוכות-טווח בריצה בודדת.

\hypertarget{ux5e7ux5e9ux5d1-ux5e8ux5d1-ux5e8ux5d0ux5e9ux5d9}{%
\subsection{קשב
רב-ראשי}\label{ux5e7ux5e9ux5d1-ux5e8ux5d1-ux5e8ux5d0ux5e9ux5d9}}

במקום לבצע פונקציית קשב אחת במלוא ממדיות המודל, ארכיטקטורת ה-Transformer
משתמשת בקשב רב-ראשי (Multi-Head Attention). הרעיון הבסיסי הוא להריץ מספר
פונקציות קשב עצמי במקביל, כל אחת עם מטריצות הטלה משלה, ולאחר מכן לאחד את
תוצאותיהן.

לכל ``ראש'' h\_i מוגדרות מטריצות הטלה פרטיות W\_i\^{}Q, W\_i\^{}K,
W\_i\^{}V ממימד d\_model x d\_k, ומטריצת פלט W\_i\^{}O ממימד d\_v x
d\_model. הקשב של ראש i מחושב כ:

head\_i = Attention(Q * W\_i\^{}Q, K * W\_i\^{}K, V * W\_i\^{}V)

ואחד הנוסחאות הפורמליות של קשב רב-ראשי היא:

MultiHead(Q, K, V) = Concat(head\_1, \ldots, head\_h) * W\^{}O

כאשר W\^{}O היא מטריצת הטלה פלט הניתנת ללמידה. בניסויים של Vaswani et
al.~(2017) נעשה שימוש ב-h=8 ראשים, עם d\_k = d\_v = d\_model / h = 64.

ההנמקה לשימוש בריבוי ראשים היא עמוקה: לכל ראש קשב ניתנת האפשרות להתמחות
בהיבט אחר של היחסים בין אסימונים. ראש אחד עשוי לתפוס תלויות תחביריות,
אחר יחסים סמנטיים, ושלישי הצמדות בין ישויות בשפה. ניסויים אמפיריים הראו
כי ראשים שונים אכן מתמחים בדפוסי קשב שונים, כגון קשרים בין שם עצם לתואר,
או קשרי אנפורה בין כינויי גוף לישויות. ריכוז תוצאותיהם של כל הראשים
באמצעות שרשור ומטריצת ההטלה W\^{}O מאפשר לדגם לנצל תובנות מרובות אלו
בו-זמנית.

\hypertarget{ux5e7ux5d9ux5d3ux5d5ux5d3-ux5deux5d9ux5e7ux5d5ux5deux5d9}{%
\subsection{קידוד
מיקומי}\label{ux5e7ux5d9ux5d3ux5d5ux5d3-ux5deux5d9ux5e7ux5d5ux5deux5d9}}

מנגנון הקשב העצמי הוא אינווריאנטי לסדר --- הוא אינו מבחין בין רצף ``כלב
נשך את האיש'' לבין ``האיש נשך את הכלב'' אם וקטורי האסימונים זהים. כדי
להזרים מידע על סדר ומיקום האסימונים ברצף, ארכיטקטורת ה-Transformer
מוסיפה לוקטורי ההטמעה (Embedding) של האסימונים קידודי מיקום (Positional
Encoding).

Vaswani et al.~(2017) הציעו קידודי מיקום סינוסואידליים, שאינם ניתנים
ללמידה אלא נקבעים מראש על פי הנוסחאות:

PE(pos, 2i) = sin(pos / 10000\^{}(2i / d\_model)) PE(pos, 2i+1) =
cos(pos / 10000\^{}(2i / d\_model))

כאשר pos הוא מיקום האסימון ברצף ו-i הוא אינדקס ממד הוקטור. כל ממד של
קידוד המיקום מתאים לגל סינוס (או קוסינוס) עם אורך גל שונה, הנע
מ-2\emph{pi לממדים קטנים ועד ל-10000}2*pi לממדים גדולים. בחירה זו מאפשרת
לקידוד לייצג מיקומים שלא נצפו באימון, ולהביע יחסי מיקום יחסי בין
אסימונים.

יתרון מכריע של הקידוד הסינוסואידלי הוא שהוא מאפשר למודל ללמוד להישען על
הזזות יחסיות: עבור הזזה קבועה k, ניתן לבטא PE(pos+k) כשילוב לינארי של
PE(pos), מה שמקל על המודל להסיק יחסי מיקום יחסי. חלופה שנחקרה באותה
עבודה היא קידוד מיקום נלמד --- מטריצת הטמעות מיקום הניתנת לאימון ---
שהביאה לתוצאות דומות, אך לוותה בהגבלה על אורך רצף מרבי.

קידודי מיקום נלמדים אומצו לאחר מכן על ידי מודלים כגון BERT ו-GPT, שבהם
רצפי האימון מסופקים תמיד בתוך חלון הקשר מוגדר מראש. מחקרים מאוחרים יותר
הציעו גישות קידוד מיקום שונות, בהן קידוד מיקום יחסי (Relative Positional
Encoding) ו-RoPE (Rotary Position Embedding), אשר שיפרו את יכולת ההכללה
לרצפים ארוכים מאלו שנראו בשלב האימון.
